
Having now introduced condensed sets -- the foundational object of condensed mathematics -- we can start to build algebra on top of them. For this we will first discuss the general theory of abelian sheaves, sheaves of rings and sheaves of modules over such rings in sections \ref{section: abelian sheaves}, \ref{section: rings, ringed spaces and ringed sites} and \ref{section: modules under morphisms}. Then afterwards, we will start to build the theory of condensed algebra in section \ref{section: condensed algebraic structures}. We will study the surrounding derived theory in chapter \ref{chapter: cohomology and derived condensed algebra}.

To construct a good theory of sheaves of modules one usually has to restrict to essentially small sites, as sheafification is a fundamental part of the theory and is only guaranteed to exist in that case. However almost all results stated solely for essentially small sites hold more generally. As long as one restricts attention to a suitable subcategory of presheaves (and the corresponding category of sheaves with such underlying presheaves) such that sheafification (as the adjoint to the forgetful functor) is defined and commutes with all limits, then virtually all the desired results hold. This is then readily satisfied by small presheaves $\smallPresheaves{\extremallyDisconnectedSets}{\Set}$ and their corresponding sheaves $\condensedSets = \smallSheaves{\extremallyDisconnectedSets}{\Set}$, where the forgetful functor admits such a left adjoint by lemma \ref{lemma: sheafification of small presheaves}. For each claim about essentially small sites, we will give a short remark about the translation to small sheaves on $\extremallyDisconnectedSets$ -- this setting is of course our main application as we wish to \emph{algebra} with condensed sets.

\section{Abelian Sheaves}
\label{section: abelian sheaves}

An important stepping stone are \emph{abelian (pre)sheaves}. Recall that there are two common ways of defining them.
\begin{definition}[Abelian (pre)sheaves]
	Let $\site{C}$ be a site. An \emph{abelian presheaf} (resp. \emph{abelian sheaf}) on $\site{C}$ is an abelian group objects internal to the category $\Presheaves{\site{C}}{\Set}$ (resp. $\Sheaves{\site{C}}{\Set}$). Equivalently, an \emph{abelian (pre)sheaf} is a (pre)sheaf of abelian groups on $\site{C}$.
	
	A \emph{morphism of abelian (pre)sheaves} is a morphism of abelian group objects. Equivalently, this is a morphism of (pre)sheaves of abelian sheaves.
\end{definition}

Of course sheafification of sheaves of sets translates to the abelian setting.
\begin{lemma}[Sheafification of abelian presheaves]
	\label{lemma: sheafification of abelian presheaves}
	
	Let $\site{C}$ be an essentially small site. Sheafification of sheaves of sets gives rise to a left adjoint $\sheafification{-}\colon \Presheaves{\site{C}}{\Ab}\to\Sheaves{\site{C}}{\Ab}$ of the forgetful functor $\Sheaves{\site{C}}{\Ab} \to \Presheaves{\site{C}}{\Ab}$ that commutes with finite limits. This left adjoint is called \emph{sheafification (of abelian presheaves)}.
\end{lemma}
\begin{proof}
	This easily seen if one considers abelian sheaves to be abelian group objects internal to sheaves of sets. The claim then follows from sheafification commuting with finite limits as per definition \ref{definition: sheafification}.
\end{proof}

\begin{remark}[Sheafification of small abelian presheaves]
	\label{remark: sheafification of small abelian presheaves}
	
	Lemma \ref{lemma: sheafification of abelian presheaves} applies verbatim to small (pre)sheaves of abelian groups on the site $\extremallyDisconnectedSets$, as sheafification $\sheafification{-}\colon \smallPresheaves{\extremallyDisconnectedSets}{\Set} \to \smallSheaves{\extremallyDisconnectedSets}{\Set}$ exists and commutes with finite limits by lemma \ref{lemma: sheafification of small presheaves}.
\end{remark}

\begin{proposition}[Abelian sheaves form an abelian category]
	\label{proposition: abelian sheaves form an abelian category}
	
	Let $\site{C}$ be an essentially small site. The category $\Sheaves{\site{C}}{\Ab}$ of abelian sheaves is abelian.
\end{proposition}
\begin{proof}
	This is a direct application \cite[\href{https://stacks.math.columbia.edu/tag/03A3}{Tag 03A3}]{stacks-project} after observing that sheafification of abelian presheaves commutes with finite limits by lemma \ref{lemma: sheafification of abelian presheaves}.
\end{proof}

\begin{remark}[Small abelian sheaves form an abelian category]
	Proposition \ref{proposition: abelian sheaves form an abelian category} applies verbatim to small abelian sheaves on the site $\extremallyDisconnectedSets$. Indeed, \cite[\href{https://stacks.math.columbia.edu/tag/03A3}{Tag 03A3}]{stacks-project} used in the proof of proposition \ref{proposition: abelian sheaves form an abelian category} applies as $\smallPresheaves{\extremallyDisconnectedSets}{\Ab}$ is abelian (all $\smallPresheaves{\boundedExtremallyDisconnectedSets{\kappa}}{\Ab}$ are abelian and the transition functors between different cardinal-cutoffs preserve the relevant structure by remark \ref{remark: transitioning in general}) and sheafification $\sheafification{-}\colon \smallPresheaves{\extremallyDisconnectedSets}{\Ab}\to \smallSheaves{\extremallyDisconnectedSets}{\Set}$ exists and is exact by remark \ref{remark: sheafification of small abelian presheaves}. Hence $\smallSheaves{\extremallyDisconnectedSets}{\Ab}$ is abelian as well.
\end{remark}

\section{Rings, Ringed Spaces \& Ringed Sites}
\label{section: rings, ringed spaces and ringed sites}

Recall the notion of ringed sites, the best-known instances of which probably being schemes.
\begin{definition}[Ringed spaces, {
		\cite[\href{https://stacks.math.columbia.edu/tag/01HA}{Tag 01HA}]{stacks-project}}]
	
	A \emph{ringed space} is a pair $(X, \struct X)$ where $X$ is a topological space and $\struct X$ is a sheaf of commutative unitary rings on $X$, called the \emph{structure sheaf}. If additionally $\stalkStruct X x$ is a local ring for any $x\in X$, then $(X, \struct X)$ is called a \emph{locally ringed space}.
	
	Let $(X,\struct X)$ and $(Y, \struct Y)$ be ringed spaces. A \emph{morphism of ringed spaces} $(X,\struct X) \to (Y, \struct Y)$ is a pair $(f, f^\sharp)$ where $f\colon X\to Y$ is continuous and $f^\sharp \colon \struct Y \to f_\ast \struct X$ is a morphism of sheaves of rings. If $(X, \struct X)$ and $(Y, \struct Y)$ are locally ringed spaces we call the pair $(f, f^\sharp)$ a \emph{morphism of locally ringed spaces} if additionally for any $x\in X$ the induced map
	$$\stalkStruct Y {f(x)} \xrightarrow{f^\sharp_{f(x)}} (f_\ast \struct X)_{f(x)} \to \stalkStruct X x$$
	on stalks is a \emph{local morphism}, \ie maps the maximal ideal to the maximal ideal.
%	We denote by $\ringedSpaces$ and $\locallyRingedSpaces$ the categories of ringed and locally ringed spaces respectively.
\end{definition}

There is a straightforward generalization (once one is comfortable replacing spaces by sites), that will prove useful when discussing condensed rings and modules over such a ring. Instead of considering sheaves on spaces, we simply generalize to sheaves on sites. First one needs a notion of sheaf of rings.
\begin{definition}[Sheaves of rings]
	Let $\site{C}$ be a site. A \emph{sheaf of rings} or simply a \emph{ring} on $\site{C}$ is a (commutative unitary) ring object internal to the category $\Sheaves{\site{C}}{\Ab}$. Equivalently, a \emph{sheaf of rings} is a sheaf of (commutative unitary) rings on $\site{C}$.
	
	A \emph{morphism of sheaves of rings} is then a morphism of ring objects. Equivalently, it is a morphism in $\Sheaves{\site{C}}{\commutativeRings}$.
\end{definition}

\begin{example}[Many topological rings give rise to sheaves of rings]
	Every topological ring $A$ whose underlying space is $\tOne$ gives rise to a sheaf of rings $\associated{A}$ (where addition and multiplication of each $\associated{A}(S)=\eHom_\Top(S, A)$ is pointwise) on the site $\extremallyDisconnectedSets$ by theorem \ref{theorem: some correspondences}. This of course includes all discrete rings, that is \quote{rings without a topology}, but also topological rings like the $p$-adic integers $\Z_p$ or more generally the ring of formal power series $\series{\Z}{T}$ with its $\ideal{T}$-adic topology.
\end{example}

With this we can provide the needed generalization of ringed spaces.
\begin{definition}[Ringed sites, {
		\cite[\href{https://stacks.math.columbia.edu/tag/03AD}{Tag 03AD}]{stacks-project}}]
	
	A \emph{ringed site} is a pair $(\site C, \structsymbol)$ where $\site C$ is a site and $\structsymbol$ is a sheaf of rings on $\site C$, called the \emph{structure sheaf}.
	
	Let $(\site C,\structsymbol)$ and $(\site C', \structsymbol')$ be ringed sites. A \emph{morphism of ringed sites} $(\site C,\structsymbol) \to (\site C', \structsymbol')$ is a pair $(f, f^\sharp)$ where $f\colon \site C\to \site C'$ is a map of sites and $f^\sharp \colon \structsymbol' \to f_\ast \structsymbol$ is a morphism of sheaves of rings, where $f_\ast \structsymbol$ carries the obvious ring structure.
%	
%	We denote by $\ringedSites$ the categories of ringed sites.
\end{definition}

Of course ringed spaces can naturally be viewed as ringed sites.
\begin{remark}[Ringed spaces as ringed sites]
	Any ringed space $(X,\struct X)$ induces a ringed site $(\open{X}, \struct X)$. Similarly, for any morphism of ringed spaces $(f,f^\sharp)\colon (X,\struct X)\to(Y,\struct Y)$, the continuous map $f\colon X\to Y$ induces a morphism of sites $\open X\to\open Y$ and $f_\ast$ (as known from the theory of sheaves on a topological space) agrees with the direct image associated to this morphism. Hence, this morphism of sites together with the map $f^\sharp$ constitutes a morphism of ringed sites $(\open X, \struct X) \to (\open Y, \struct Y)$.
%	We obtain a faithful functor
%		$$\ringedSpaces \to \ringedSites.$$
\end{remark}

\begin{remark}[Open subspaces of ringed spaces]
	Let $(X,\struct X)$ be a ringed space and $U\subseteq X$ open. Then $\struct U\ldef \struct X\restriction_U$ is a sheaf of rings on $U$ and hence $(U, \struct U)$ naturally is a ringed space. Even more, there is a natural map of ringed spaces $(i,i^\sharp)\colon (U, \struct U)\to (X, \struct X)$ where $i\colon U\to X$ is the inclusion and $i^\sharp_V \colon \struct X(V)\to (i_\ast\struct U)(V) = \struct X(U\cap V)$ is the restriction map for all $V$ open in $X$.
\end{remark}


\section{(Pre)sheaves of Modules}
\label{section: (pre)sheaves of modules}

\begin{definition}[(Pre)sheaves of modules, {
	\cite[\href{https://stacks.math.columbia.edu/tag/03CV}{Tag 03CV}]{stacks-project}
	\&
	\cite[\href{https://stacks.math.columbia.edu/tag/03CT}{Tag 03CT}]{stacks-project}}]
	\label{definition: (pre)sheaves of modules}
	
	Let $(\site C, \structsymbol)$ be a ringed site. A \emph{presheaf of $\structsymbol$-modules} $\sheaf F$ is a module object over the sheaf of rings $\structsymbol$, internal to the category $\Presheaves{\site C}{\Ab}$. A \emph{sheaf of $\structsymbol$-modules} is a presheaf of $\structsymbol$-modules whose underlying presheaf is a sheaf.
	
	A \emph{morphism of presheaves of $\structsymbol$-modules} $\sheaf F\to\sheaf G$ is a morphism of the corresponding module objects internal to $\Presheaves{\site C}{\Ab}$. A \emph{morphism of sheaves of $\structsymbol$-modules} is a morphism of the corresponding presheaves of $\structsymbol$-modules.
	
	We denote by $\pMod{\structsymbol}$ and $\mod\structsymbol$ the categories of presheaves and sheaves of $\structsymbol$-modules respectively. Sets of morphisms in $\pMod\structsymbol$ and $\mod\structsymbol$ are denoted by $\Hom_\structsymbol$.
\end{definition}

As definition \ref{definition: (pre)sheaves of modules} is quite abstract, it makes sense to try to gain a more explicit understanding of (pre)sheaves of modules and their morphisms. We will do so by inspecting them on the level of sections.
\begin{remark}[(Pre)sheaves of modules, more explicit]
	The condition in definition \ref{definition: (pre)sheaves of modules} for a sheaf $\sheaf{F}$ to be a (pre)sheaf of $\structsymbol$-modules can be given more explicitly. Indeed, we require, that for any $U \in \site C$ the sections $\sheaf F(U)$ carry the structure of an $\structsymbol(U)$-module and for any morphism $V\subseteq U$ the restriction $\sheaf F(U)\to \sheaf F(V)$ is $\structsymbol(U)$-linear where $\sheaf F(V)$ is considered as an $\structsymbol(U)$-module via restriction along the morphism $\structsymbol(U)\to\structsymbol(V)$. Concretely, for any section $s\in \sheaf F(U)$ and any $r\in \structsymbol(U)$ we require $(r\cdot s)\restriction_V = r\restriction_V \cdot s\restriction_V$.
	
	Similarly, the condition on being a morphism of (pre)sheaves of $\structsymbol$-modules can be rephrased as follows: An additive morphism $f\colon \sheaf F\to \sheaf G$ between (pre)sheaves of $\structsymbol$-modules is a morphism of (pre)sheaves of $\structsymbol$-modules if
	\begin{center}
		\begin{tikzcd}
			\structsymbol \times \sheaf F &\sheaf F\\
			\structsymbol \times \sheaf G &\sheaf G
			\ar[from=1-1, to=1-2]
			\ar[from=2-1, to=2-2]
			\ar[from=1-1, to=2-1, "\id \times f", swap]
			\ar[from=1-2, to=2-2, "f"]
		\end{tikzcd}
	\end{center}
	commutes, \ie for any $U\subseteq X$, any $r\in \structsymbol(U)$ and any $s\in \sheaf F(U)$ we require $f_U(r\cdot s) = r\cdot f_U(s)$.
\end{remark}

As one expects, there is a canonical way of turning a presheaf of modules into a sheaf of modules. This is formalized in the following lemma.
\begin{lemma}[Sheafification of presheaves of modules]
	\label{lemma: sheafification of presheaves of modules}
	
	Let $(\site C, \structsymbol)$ be an essentially small ringed site. For each presheaf of $\structsymbol$-modules $\sheaf F$ there is a unique $\structsymbol$-module structure $\structsymbol\times \sheaf F^\sheafificationsymbol \to \sheaf F^\sheafificationsymbol$ on $\sheaf F^\sheafificationsymbol$ such that
	\begin{center}
		\begin{tikzcd}
			\structsymbol \times \sheaf F &\sheaf F\\
			\structsymbol \times \sheaf F^\sheafificationsymbol & \sheaf F^\sheafificationsymbol
			\ar[from=1-1, to=1-2]
			\ar[from=1-1, to=2-1]
			\ar[from=1-2, to=2-2]
			\ar[from=2-1, to=2-2]
		\end{tikzcd}
	\end{center}
	commutes. Additionally, if $\sheaf G$ is a sheaf of $\structsymbol$-modules, then any morphism $\sheaf F\to\sheaf G$ of presheaves of $\structsymbol$-modules factors uniquely as $\sheaf F\to \sheaf F^\sheafificationsymbol\to\sheaf G$ where $\sheaf F^\sheafificationsymbol \to \sheaf G$ is a morphism of sheaves of $\structsymbol$-modules. Hence sheafification of abelian sheaves gives rise to an adjunction
		$$\mod\structsymbol \xleftrightarrows{\sheafification{-}}{} \pMod\structsymbol.$$
\end{lemma}
\begin{proof}
	This follows from \cite[\href{https://stacks.math.columbia.edu/tag/03CY}{Tag 03CY}]{stacks-project} using that the unit $\structsymbol\to\structsymbol ^\sheafificationsymbol$ is an isomorphism.
\end{proof}

\begin{remark}[Sheafification of small presheaves of modules]
	The previous lemma \ref{lemma: sheafification of presheaves of modules} applies the small presheaves of modules on the site $\extremallyDisconnectedSets$ for any small sheaf of rings. Indeed, this is a direct consequence of the existence of sheafification and its commutation with finite limits.
\end{remark}

We also find that the category of sheaves of modules are abelian.
\begin{proposition}[The abelian category of sheaves of modules]
	\label{proposition: the abelian category of sheaves of modules}
	
	Let $(\site C, \structsymbol)$ be an essentially small ringed site. The category $\mod\structsymbol$ of sheaves of $\structsymbol$-modules is an abelian category. The forgetful functor $\mod\structsymbol \to \Sheaves{\site C}{\Ab}$ is exact, hence kernels, cokernels and exactness of sequences of sheaves of $\structsymbol$-modules agree with their corresponding notions of abelian sheaves.
\end{proposition}
\begin{shortproof}
	This is \cite[\href{https://stacks.math.columbia.edu/tag/03DA}{Tag 03DA}]{stacks-project}.
\end{shortproof}

\begin{remark}
	While proposition \ref{proposition: the abelian category of sheaves of modules} translates to small sheaves as well, it is more instructive to state the corresponding result later on in theorem \ref{theorem: condensed modules form an abelian category} as there is much more to say about such categories. 
\end{remark}

\begin{lemma}[Tensor products \& internal $\Hom$s]
	\label{lemma: tensor products and internal Homs}
	
	Let $(\site{C}, \structsymbol)$ be a ringed site. Then $\pMod{\structsymbol}$ admits a symmetric monoidal tensor product given on presheaves $\sheaf{F}$ and $\sheaf{G}$ by
		$$\sheaf F \otimes_{p, \structsymbol} \sheaf{G} \ldef \mleft(U\mapsto \sheaf{F}(U)\otimes_{\structsymbol(U)} \sheaf{G}(U)\mright).$$
	This tensor product admits a partial adjoint given on presheaves $\sheaf{F}$ and $\sheaf{G}$ by
		$$\intHom_\structsymbol(\sheaf{F}, \sheaf{G})\ldef \mleft(U\mapsto \Hom_\structsymbol(\sheaf{F}\resTo{U}, \sheaf{G}\resTo{U})\mright),$$
	where $\sheaf{F}\resTo{U}$ and $\sheaf{G}\resTo{U}$ are the induced presheaves on the slice-category $\site{C}/U$. This makes $\pMod{\structsymbol}$ a closed symmetric monoidal category.
	
	If $\site{C}$ is essentially small, the category $\mod{\structsymbol}$ of sheaves of $\structsymbol$-modules is closed symmetric monoidal as well, with tensor product $-\otimes_\structsymbol- \ldef \sheafification{-\otimes_{p, \structsymbol}-}$ and internal $\Hom$ given by $\intHom_\structsymbol$.
\end{lemma}
\begin{proof}
	The existence of the tensor product and the internal Hom (as well as their properties)
	can be found at \cite[\href{https://stacks.math.columbia.edu/tag/03EK}{Tag 03EK}]{stacks-project} and \cite[\href{https://stacks.math.columbia.edu/tag/04TT}{Tag 04TT}]{stacks-project} respectively. For the claimed adjunctions consider \cite[\href{https://stacks.math.columbia.edu/tag/03EO}{Tag 03EO}]{stacks-project}.
\end{proof}

The partial adjunctions even enrich to the internal $\Hom$.
\begin{lemma}[(Enriched) tensor-Hom adjunction]
	\label{lemma: enriched tensor-Hom adjunction}
	
	Let $(\site C, \structsymbol)$ be an essentially small ringed site. There are isomorphisms
	$$\intHom_\structsymbol(\sheaf F\otimes_\structsymbol\sheaf G, \sheaf H) \cong \intHom_\structsymbol(\sheaf F, \intHom_\structsymbol(\sheaf G, \sheaf H))$$
	natural in the sheaves $\sheaf F$, $\sheaf G$ and $\sheaf H$ of $\structsymbol$-modules.
\end{lemma}
\begin{shortproof}
	This is \cite[\href{https://stacks.math.columbia.edu/tag/03EO}{Tag 03EO}]{stacks-project}.
\end{shortproof}

\begin{remark}[Tensor products and internal $\Hom$s of small sheaves]
	Again, by the usual argument, lemma \ref{lemma: tensor products and internal Homs} and lemma \ref{lemma: enriched tensor-Hom adjunction} translate to small (pre)sheaves of modules on $\extremallyDisconnectedSets$ for any small sheaf of rings. 
\end{remark}

\section{Modules under Morphisms}
\label{section: modules under morphisms}

\begin{lemma}[Direct and inverse images of (pre)sheaves of modules]
	\label{lemma: direct and inverse images of (pre)sheaves of modules}
	
	Let $f\colon \site C\to\site C'$ be a morphism of essentially small sites and let $\structsymbol$ and $\structsymbol'$ be sheaves of rings on $\site C$ and $\site C'$ respectively. Let $\sheaf F$ be a (pre)sheaf of $\structsymbol$-modules. There is a natural morphism $f_\ast\structsymbol \times f_\ast \sheaf F\to f_\ast\sheaf F$ that turns $f_\ast\sheaf F$ into a (pre)sheaf of $f_\ast\structsymbol$-modules. Similarly, for any (pre)sheaf $\sheaf G$ of $\structsymbol'$-modules there is a natural morphism $f^{-1}\structsymbol' \times f^{-1}\sheaf G\to f^{-1}\sheaf G$ which turns $f^{-1}\sheaf G$ into a (pre)sheaf of $f^{-1}\structsymbol'$-modules. Both constructions are functorial, that is we obtain functors
		$$\maybePMod\structsymbol \to \maybePMod{f_\ast\structsymbol}$$
	and
		$$\maybePMod{\structsymbol'}\to\maybePMod{f^{-1}\structsymbol'}.$$
\end{lemma}
\begin{proof}
	This is shown in \cite[\href{https://stacks.math.columbia.edu/tag/03D1}{Tag 03D1}]{stacks-project} and \cite[\href{https://stacks.math.columbia.edu/tag/03D2}{Tag 03D2}]{stacks-project} and boils down to the fact that both $f_\ast$ and $f^{-1}$ commute with finite products of sheaves of sets.
\end{proof}

\begin{remark}[Direct and inverse images of (pre)sheaves of modules]
	\label{remark: direct and inverse images of (pre)sheaves of modules}
	
	One has to be careful when translating lemma \ref{lemma: direct and inverse images of (pre)sheaves of modules} to small sheaves. First, it is necessary to distinguish the cases of morphisms of sites $\extremallyDisconnectedSets\to\site{C'}$ and $\site{C} \to \extremallyDisconnectedSets$ and even endomorphisms $\extremallyDisconnectedSets \to \extremallyDisconnectedSets$ (essentially by virtue that we defined smallness only for the site $\extremallyDisconnectedSets$, not in general). Secondly, it is not clear the usual definitions of $f_\ast$ and $f^{-1}$ simply translate. In an exchange with \href{https://home.sandiego.edu/~shulman/}{Professor Mike Schulman}, he mentioned that in general (contrary to what one might initially expect) a morphism of sites always induces an inverse image $f^{-1}$ of small sheaves (small in the sense of his paper \cite{shulman2012exactcompletionssmallsheaves}, \confer remarks \ref{remark: shulman's small sheaves and independence of the defining site} and \ref{remark: conflicting terminology}), but what fails is the existence of the direct image! The author is however not aware of any reference of this fact, as the literature on small sheaves (again, in the sense of Shulman) is rather thin. He could not locate the result in \eg \cite{shulman2012exactcompletionssmallsheaves}.
	
	Either way, in this thesis we will only need the very special case with $f^\top = \id_\extremallyDisconnectedSets$ being the identity. In this case of course $f_\ast$ and $f^{-1}$ are simply the identity and the result follows. This will prove useful in defining scalar restriction and extension of modules on $\extremallyDisconnectedSets$ later on.
\end{remark}

\begin{lemma}[Pullback and pushforward of sheaves of modules, {
	\cite[\href{https://stacks.math.columbia.edu/tag/03D6}{Tag 03D6}]{stacks-project}}]
	\label{lemma: pullback and pushforward of sheaves of modules}
	
	Let $(f,f^\sharp)\colon (\site C, \structsymbol) \to (\site C',\structsymbol')$ be a morphism of essentially small ringed sites. We define two functors
		$$f_\ast\colon \mod\structsymbol\rightarrow\mod{\structsymbol'}$$
	and
		$$f^\ast\colon \mod\structsymbol\leftarrow \mod{\structsymbol'}.$$
	\begin{itemize}
		\item
		Let $\sheaf F$ be a sheaf of $\structsymbol$-modules. The \emph{pushforward of $\sheaf F$ along $f$} is the sheaf of $\structsymbol'$-modules with underlying abelian sheaf $f_\ast \sheaf F$ and module structure given by the map $f^\sharp\colon \structsymbol'\to f_\ast \structsymbol$ and the $f_\ast\structsymbol$-module structure $f_\ast\structsymbol\times f_\ast\sheaf F\to f_\ast\sheaf F$ of lemma \ref{lemma: direct and inverse images of (pre)sheaves of modules}.
		
		\item
		Let $\sheaf G$ be a sheaf of $\structsymbol'$-modules. The \emph{pullback of $\sheaf G$ along $f$} is the sheaf of $\structsymbol$-modules with underlying abelian sheaf
			$$f^\ast\sheaf G\ldef \structsymbol\otimes_{f^{-1}\structsymbol'} f^{-1}\sheaf G$$
		where $\structsymbol$ is a sheaf of $f^{-1}\structsymbol'$-modules by virtue of the map $f^{-1}\structsymbol'\to\structsymbol$ corresponding to $f^\sharp\colon \structsymbol'\to f_\ast\structsymbol$ under the adjunction $f^{-1}\ladj f_\ast$. The module structure again is induced by lemma \ref{lemma: direct and inverse images of (pre)sheaves of modules} after applying $\structsymbol \otimes_{f^{-1}\structsymbol'} -$ and using the isomorphism $\structsymbol \to \structsymbol\otimes_{f^{-1}\structsymbol'} f^{-1}\structsymbol$.
	\end{itemize}
	
	By \cite[\href{https://stacks.math.columbia.edu/tag/03D7}{Tag 03D7}]{stacks-project}, the pushforward and pullback form and adjunction
		$$\mod{\structsymbol}\xleftrightarrows{f^\ast}{f_\ast}\mod{\structsymbol'},$$
	that is, the pullback $f^\ast$ along $f$ is left adjoint to the pushforward $f_\ast$ along $f$.
\end{lemma}

\begin{corollary}[Change of rings]
	\label{corollary: change of rings}
	
	Let $\site C$ be an essentially small site and $\phi\colon \structsymbol \to \structsymbol'$ be a morphism of rings on $\site C$. There is an adjoint pair of functors
		$$\mod {\structsymbol'} \xleftrightarrows{\scalarExtension{\phi}}{\scalarRestriction{\phi}} \mod\structsymbol$$
	where $\scalarExtension{\phi}\sheaf G \ldef \structsymbol' \otimes_\structsymbol \sheaf G$. The functors $\scalarExtension{\phi}$ and $\scalarRestriction{\phi}$ are called \emph{scalar extension} and \emph{scalar restriction} respectively.
\end{corollary}
\begin{proof}
	This trivially follows from lemma \ref{lemma: pullback and pushforward of sheaves of modules} observing that $(\id_\site C, \phi)\colon (\site C,\structsymbol')\to(\site C, \structsymbol)$ is a morphism of ringed sites.
\end{proof}


\begin{remark}[Pullback and pushforward of small sheaves]
	As in remark \ref{remark: direct and inverse images of (pre)sheaves of modules}, for translating lemma \ref{lemma: pullback and pushforward of sheaves of modules} to small sheaves, we restrict to the special case of $f^\top = \id_\extremallyDisconnectedSets$ being the identity. In this case the proof immediately follows from \cite[\href{https://stacks.math.columbia.edu/tag/03CZ}{Tag 03CZ}]{stacks-project}. This is then also precisely corollary \ref{corollary: change of rings} for the site $\extremallyDisconnectedSets$.
\end{remark}

It will prove useful later on to know that scalar restriction is always exact.
\begin{lemma}[Scalar restriction is exact]
	Let $\site{C}$ be an essentially small site and $\phi\colon \structsymbol \to \structsymbol'$ be a morphism of rings on $\site{C}$. Then scalar restriction $\scalarRestriction{\phi}\colon \mod{\structsymbol'}\to\mod{\structsymbol}$ is exact.
\end{lemma}
\begin{proof}
	As a right adjoint, $\scalarRestriction{\phi}$ certainly is left exact. It remains to show that $\scalarRestriction{\phi}$ is right exact. Recall from proposition \ref{proposition: the abelian category of sheaves of modules} that both categories of modules are abelian and that exactness can be verified on the level of abelian sheaves. The claim then follows from \cite[\href{https://stacks.math.columbia.edu/tag/04DB}{Tag 04DB}]{stacks-project} after observing that the geometric morphism of topoi $\Sheaves{\site{C}}{\Set} \xleftrightarrows{}{} \Sheaves{\site{C}}{\Set}$ induced by $\id_\site{C}\colon \site{C}\to \site{C}$ trivially satisfies the four equivalent conditions. 
\end{proof}

\begin{lemma}[Free sheaves]
	\label{lemma: free sheaves}
	
	Let $\site C$ be an essentially small site. 
	\begin{enumerate}
		\item
		The forgetful functor $\Sheaves{\site C}{\Ab} \to \Sheaves{\site C}{\Set}$ admits a left adjoint $\free {\constantSheaf{\Z}} -$, such that
			$$\free {\constantSheaf{\Z}} T = \sheafification{U \mapsto \free \Z {T(U)}},$$
		for any sheaf of sets $T$. We call $\free {\constantSheaf{\Z}} T$ the \emph{free abelian group on $T$}.
		
		\item
		Let $\structsymbol$ be a sheaf of rings on $\site C$. The forgetful functor $\mod\structsymbol \to \Sheaves{\site C}{\Set}$ admits a left adjoint $\free \structsymbol -$, such that
			$$\free \structsymbol T = \structsymbol \otimes_{\constantSheaf\Z} \free {\constantSheaf \Z} T$$
		is the extension of scalars along the unique map $\constantSheaf{\Z} \to \structsymbol$, for any sheaf of sets $T$ . We call $\free \structsymbol T$ the \emph{free $\structsymbol$-module on $T$}.
	\end{enumerate}
\end{lemma}
\begin{proof}
	Consider the functor $\free{\associated{\Z}^p}{-}\colon\Presheaves{\site{C}}{\Set} \to \Presheaves{\site{C}}{\Ab}$ given by $\sheaf{F}\mapsto \mleft(U\mapsto \free{\Z}{\sheaf{F}(U)}\mright)$. Let $T$ be a presheaf of sets and $A$ be a presheaf of abelian groups. Then there is a natural morphism
		$$\Hom_\Presheaves{\site{C}}{\Ab}(\free{\associated{\Z}^p}{T}, A) \to \Hom_\Presheaves{\site{C}}{\Set}(T, A)$$
	given by mapping a natural transformation $\eta$ to the natural transformation with component $T(U)\hookrightarrow\free{\Z}{T(U)} \xrightarrow{\eta_U} A(U)$ at $U\in\site{C}$. By the universal property of $\free{\Z}{-}$ this natural morphism admits and inverse, hence $\free{\associated{\Z}^p}{-}$ is a left adjoint to the forgetful functor $\Presheaves{\site{C}}{\Ab} \to \Presheaves{\site{C}}{\Set}$. Define $\free{\associated{\Z}}{-}\ldef \free{\associated{\Z}^p}{-}^\sheafificationsymbol\colon \Sheaves{\site{C}}{\Set}\to\Sheaves{\site{C}}{\Ab}$. Now if $T$ is a sheaf of sets and $A$ is an abelian sheaf, then
		$$\Hom_\sheavesSymbol(\free{\associated{\Z}}{T}, A) \cong \Hom_\preSheavesSymbol(\free{\associated{\Z}^p}{T}, A) \cong \Hom_\preSheavesSymbol(T, A)=\Hom_\sheavesSymbol(T, A)$$
	showing that $\free{\associated{\Z}}{-}$ is the desired adjoint. Denote by $\phi$ the unique morphism of rings $\associated{\Z}\to \structsymbol$. Since the forgetful functor from sheaves of $\structsymbol$-modules to sheaves of sets factors over the forgetful functor from sheaves of abelian groups to sheaves of sets, it is also clear that $\free{\structsymbol}{-} \ldef \scalarExtension{\phi}\free{\associated{\Z}}{-} = \structsymbol\otimes_{\associated{\Z}}\free{\associated{\Z}}{-}$ is left adjoint to $\mod{\structsymbol}\to\Sheaves{\site{C}}{\Set}$.
\end{proof}

\begin{remark}[Free small sheaves]
	\label{remark: free small sheaves}
	
	Lemma \ref{lemma: free sheaves} translates well to small sheaves, as the proof only relies on the existence of sheafification as an adjoint to the forgetful functor. It is also clear that for a small (pre)sheaf $\sheaf F$, the presheaf $\free{\associated{\Z}^p}{\sheaf{F}}$ is again small, as the analogous functors for $\boundedExtremallyDisconnectedSets{\kappa}$ commute with the transition functors $(-)^\enlargementsymbol{\kappa}{\kappa'}$.
\end{remark}
